\textbf{Phần 1. Trắc nghiệm lựa chọn}
\setcounterpageref{ex}{0}

\begin{ex} %[5D6V30]
Cho 4 chữ số $0, 3, 5, 8$. Số thập phân bé nhất có đủ 4 chữ số đã cho và có đúng hai chữ số ở phần thập phân là
\choice
{$03{,}58$}
{$3{,}058$}
{\True $30{,}58$}
{$35{,}08$}
\loigiai{
Để số thập phân là bé nhất thì chữ số ở hàng cao nhất của phần nguyên phải bé nhất (khác $0$), do đó phần nguyên là $30$.\\
Hai chữ số ở phần thập phân xếp theo thứ tự tăng dần là $5$ và $8$.\\
Số thập phân bé nhất cần tìm là $30{,}58$.\\
\textbf{Đáp án của câu hỏi là C}
}
\end{ex}

\begin{ex} %[5D6V30]
Tìm chữ số chẵn lớn nhất của $x$ để: $45{,}7x8 < 45{,}738$.
\choice
{$0$}
{$1$}
{\True $2$}
{$4$}
\loigiai{
So sánh hai số có cùng các hàng chục, đơn vị, phần mười và phần nghìn:\\
Để $45{,}7x8 < 45{,}738$ thì $x < 3$.\\
Vì $x$ là chữ số chẵn nên $x \in \{0; 2\}$.\\
Chữ số chẵn lớn nhất là $x = 2$.\\
\textbf{Đáp án của câu hỏi là C}
}
\end{ex}

\begin{ex} %[5D6V30]
Làm tròn số thập phân $68{,}475$ đến hàng phần mười ta được số
\choice
{$68{,}4$}
{\True $68{,}5$}
{$68{,}48$}
{$69$}
\loigiai{
Chữ số ngay sau hàng phần mười (hàng phần trăm) là chữ số $7 \ge 5$.\\
Do đó ta cộng thêm $1$ vào chữ số hàng phần mười: $4 + 1 = 5$.\\
Kết quả làm tròn là $68{,}5$.\\
\textbf{Đáp án của câu hỏi là B}
}
\end{ex}

\begin{ex} %[5D6V30]
Viết số đo $5\text{ m}^2\text{ }8\text{ dm}^2$ dưới dạng số thập phân có đơn vị đo là mét vuông ta được
\choice
{$5{,}8\text{ m}^2$}
{\True $5{,}08\text{ m}^2$}
{$5{,}008\text{ m}^2$}
{$58\text{ m}^2$}
\loigiai{
Vì $1\text{ m}^2 = 100\text{ dm}^2$ nên $8\text{ dm}^2 = 0{,}08\text{ m}^2$.\\
Vậy $5\text{ m}^2\text{ }8\text{ dm}^2 = 5{,}08\text{ m}^2$.\\
\textbf{Đáp án của câu hỏi là B}
}
\end{ex}

\begin{ex} %[5D6V30]
Dãy số thập phân nào dưới đây được sắp xếp theo thứ tự từ bé đến lớn?
\choice
{$0{,}53;\ 0{,}503;\ 0{,}35;\ 0{,}305$}
{$0{,}35;\ 0{,}305;\ 0{,}503;\ 0{,}53$}
{\True $0{,}305;\ 0{,}35;\ 0{,}503;\ 0{,}53$}
{$0{,}305;\ 0{,}503;\ 0{,}35;\ 0{,}53$}
\loigiai{
So sánh các số thập phân có phần nguyên bằng $0$:\\
$0{,}305 < 0{,}350 < 0{,}503 < 0{,}530$.\\
Vậy thứ tự từ bé đến lớn là: $0{,}305;\ 0{,}35;\ 0{,}503;\ 0{,}53$.\\
\textbf{Đáp án của câu hỏi là C}
}
\end{ex}

\begin{ex} %[5D6V30]
Khi dời dấu phẩy của một số thập phân sang bên phải một chữ số thì chữ số ở hàng phần mười của số ban đầu chuyển thành chữ số ở hàng nào của số mới?
\choice
{Hàng chục}
{\True Hàng đơn vị}
{Hàng phần trăm}
{Hàng phần mười}
\loigiai{
Khi dời dấu phẩy sang phải một chữ số, số đó gấp lên $10$ lần.\\
Hàng phần mười nhân với $10$ sẽ thành hàng đơn vị.\\
\textbf{Đáp án của câu hỏi là B}
}
\end{ex}

\begin{ex} %[5D6V30]
Có bao nhiêu số thập phân có đúng một chữ số ở phần thập phân nằm giữa hai số $4{,}5$ và $5{,}6$?
\choice
{$9\text{ số}$}
{\True $10\text{ số}$}
{$11\text{ số}$}
{Vô số}
\loigiai{
Các số thập phân có 1 chữ số ở phần thập phân lớn hơn $4{,}5$ và nhỏ hơn $5{,}6$ là:\\
$4{,}6;\ 4{,}7;\ 4{,}8;\ 4{,}9;\ 5{,}0;\ 5{,}1;\ 5{,}2;\ 5{,}3;\ 5{,}4;\ 5{,}5$.\\
Tổng cộng có $10$ số.\\
\textbf{Đáp án của câu hỏi là B}
}
\end{ex}

\begin{ex} %[5D6V30]
Giá trị của chữ số $7$ trong số $15{,}72$ gấp bao nhiêu lần giá trị của chữ số $7$ trong số $23{,}475$?
\choice
{$1\text{ lần}$}
{\True $10\text{ lần}$}
{$100\text{ lần}$}
{$1\,000\text{ lần}$}
\loigiai{
Chữ số $7$ trong số $15{,}72$ có giá trị là $0{,}7$ (hàng phần mười).\\
Chữ số $7$ trong số $23{,}475$ có giá trị là $0{,}07$ (hàng phần trăm).\\
Tỉ số là: $0{,}7 : 0{,}07 = 10$ (lần).\\
\textbf{Đáp án của câu hỏi là B}
}
\end{ex}

\begin{ex} %[5D6V30]
Tìm số tự nhiên chẵn $x$ thỏa mãn: $12{,}85 < x < 15{,}12$.
\choice
{$12$}
{\True $14$}
{$16$}
{$13$}
\loigiai{
Các số tự nhiên nằm giữa $12{,}85$ và $15{,}12$ là: $13;\ 14;\ 15$.\\
Trong các số trên, số tự nhiên chẵn là $14$.\\
\textbf{Đáp án của câu hỏi là B}
}
\end{ex}

\begin{ex} %[5D6V30]
Viết phân số thập phân $\dfrac{375}{1\,000}$ dưới dạng số thập phân ta được
\choice
{$3{,}75$}
{$37{,}5$}
{\True $0{,}375$}
{$0{,}0375$}
\loigiai{
Chia cho $1\,000$, ta dời dấu phẩy sang trái 3 chữ số: $\dfrac{375}{1\,000} = 0{,}375$.\\
\textbf{Đáp án của câu hỏi là C}
}
\end{ex}

\begin{ex} %[5D6V30]
Bác nông dân cân 4 bao thóc có khối lượng lần lượt là: $48{,}5\text{ kg};\ 48{,}75\text{ kg};\ 48{,}6\text{ kg};\ 48{,}09\text{ kg}$. Bao thóc nặng nhất là bao có khối lượng
\choice
{$48{,}5\text{ kg}$}
{\True $48{,}75\text{ kg}$}
{$48{,}6\text{ kg}$}
{$48{,}09\text{ kg}$}
\loigiai{
So sánh các số đo: $48{,}09 < 48{,}50 < 48{,}60 < 48{,}75$.\\
Bao nặng nhất là bao có khối lượng $48{,}75\text{ kg}$.\\
\textbf{Đáp án của câu hỏi là B}
}
\end{ex}

\begin{ex} %[5D6V30]
Đổi $3\text{ tấn }45\text{ kg}$ ra tấn ta được số thập phân là
\choice
{$3{,}45\text{ tấn}$}
{\True $3{,}045\text{ tấn}$}
{$3{,}0045\text{ tấn}$}
{$34{,}5\text{ tấn}$}
\loigiai{
Vì $1\text{ tấn} = 1\,000\text{ kg}$ nên $45\text{ kg} = 0{,}045\text{ tấn}$.\\
Do đó $3\text{ tấn }45\text{ kg} = 3{,}045\text{ tấn}$.\\
\textbf{Đáp án của câu hỏi là B}
}
\end{ex}

\textbf{Phần 2. Trắc nghiệm đúng sai}
\setcounterpageref{ex}{0}

\begin{ex} %[5D6V30]
Cho số thập phân $A = 125{,}678$.
\choiceTF
{\True Phần nguyên của số $A$ gồm $1$ trăm, $2$ chục và $5$ đơn vị}
{Chữ số $7$ trong số $A$ có giá trị thuộc hàng phần mười}
{\True Làm tròn số $A$ đến hàng phần trăm ta được số $125{,}68$}
{\True Làm tròn số $A$ đến hàng đơn vị ta được số $126$}
\loigiai{
\begin{itemchoice}
\item Đúng. Phần nguyên là $125$ gồm $1$ trăm, $2$ chục, $5$ đơn vị.
\item Sai. Chữ số $7$ đứng ở vị trí thứ hai sau dấu phẩy nên thuộc hàng phần trăm (giá trị là $0{,}07$).
\item Đúng. Chữ số sau hàng phần trăm là $8 \ge 5$ nên làm tròn lên $125{,}68$.
\item Đúng. Chữ số sau hàng đơn vị là $6 \ge 5$ nên làm tròn lên $126$.
\end{itemchoice}
}
\end{ex}

\begin{ex} %[5D6V30]
Xét tính đúng sai của các khẳng định chuyển đổi đơn vị và so sánh sau:
\choiceTF
{\True $4\text{ m }5\text{ cm} = 4{,}05\text{ m}$}
{Khoảng thời gian $2{,}5\text{ giờ} = 2\text{ giờ }50\text{ phút}$}
{\True $3\text{ ha }400\text{ m}^2 = 3{,}04\text{ ha}$}
{\True Số thập phân $0{,}9$ lớn hơn số thập phân $0{,}899$}
\loigiai{
\begin{itemchoice}
\item Đúng. Vì $1\text{ m} = 100\text{ cm}$ nên $5\text{ cm} = 0{,}05\text{ m}$.
\item Sai. Vì $0{,}5\text{ giờ} = 30\text{ phút}$ nên $2{,}5\text{ giờ} = 2\text{ giờ }30\text{ phút} 
eq 2\text{ giờ }50\text{ phút}$.
\item Đúng. Vì $1\text{ ha} = 10\,000\text{ m}^2$ nên $400\text{ m}^2 = 0{,}04\text{ ha}$.
\item Đúng. So sánh hàng phần mười: $9 > 8$ nên $0{,}9 > 0{,}899$.
\end{itemchoice}
}
\end{ex}

\textbf{Phần 3. Trắc nghiệm trả lời ngắn}
\setcounterpageref{ex}{0}

\begin{ex} %[5D6V30]
Có bao nhiêu số tự nhiên $x$ thỏa mãn: $8{,}4 < x < 15{,}2$?
\shortans{7}
\loigiai{
Các số tự nhiên $x$ thỏa mãn là: $9;\ 10;\ 11;\ 12;\ 13;\ 14;\ 15$.\\
Tổng cộng có 7 số tự nhiên thỏa mãn.
}
\end{ex}

\begin{ex} %[5D6V30]
Tìm số tự nhiên lớn nhất có hai chữ số mà nhỏ hơn $85{,}7$.
\shortans{85}
\loigiai{
Số tự nhiên lớn nhất nhỏ hơn $85{,}7$ là số $85$.
}
\end{ex}

\begin{ex} %[5D6V30]
Cho ba chữ số $1, 2, 4$. Viết số thập phân lớn nhất có một chữ số ở phần nguyên và hai chữ số ở phần thập phân từ ba chữ số đã cho (mỗi chữ số dùng một lần). Nhân số đó với $100$ ta được kết quả là bao nhiêu?
\shortans{421}
\loigiai{
Số thập phân lớn nhất lập được là $4{,}21$.\\
Nhân số đó với 100: $4{,}21 \times 100 = 421$.
}
\end{ex}

\begin{ex} %[5D6V30]
Khi dời dấu phẩy của một số thập phân sang bên phải một chữ số, ta được số mới lớn hơn số ban đầu là $117$ đơn vị. Tìm số thập phân ban đầu.
\shortans{13}
\loigiai{
Số mới gấp 10 lần số ban đầu, tức là hơn số ban đầu 9 lần.\\
Số thập phân ban đầu là:
$117 : 9 = 13$.
}
\end{ex}

\begin{ex} %[5D6V30]
Có bao nhiêu số thập phân có đúng một chữ số ở phần thập phân thỏa mãn điều kiện: $2 < x < 5$?
\shortans{29}
\loigiai{
Các số đó là: $2{,}1;\ 2{,}2;\ \dots;\ 4{,}9$.\\
Số lượng các số là:
$(4{,}9 - 2{,}1) : 0{,}1 + 1 = 29\text{ (số)}$.
}
\end{ex}

\begin{ex} %[5D6V30]
Một cuộn vải dài $35\text{ m}$. Người ta cắt thành các đoạn nhỏ bằng nhau, mỗi đoạn dài $2{,}5\text{ m}$. Hỏi cắt được tất cả bao nhiêu đoạn vải như thế?
\shortans{14}
\loigiai{
Số đoạn vải cắt được là:
$35 : 2{,}5 = 14\text{ (đoạn)}$.
}
\end{ex}
